\scanpage{67}
Passage du splitting standard des formes différentielles de 2\textsuperscript{\unsure{ème|e}} espèce en $\omega,\eta$ au splitting \unsure{transcendant} (de Hodge) $\omega,\bar\omega$, donc on aura
\[
\left\{\begin{aligned}
\eta+\alpha\omega&=\beta\bar{\omega}\\
\alpha&=-\frac{\bar{\omega}_1\eta_2-\bar{\omega}_2\eta_1}
{\bar{\omega}_1\omega_2-\bar{\omega}_2\omega_1}
\end{aligned}\right.
\]
\margin{N.B. dénominateur $=2i\,|\omega_1|^2\,\mathrm{Im}\,\tau\ne0$}
\uline{NB}\quad On a une 2 forme \unsure{invariante|invariant} canonique associée à la comparaison des deux splittings
\[
\pi=\alpha\omega^2\quad\in\unsure{\Gamma}\operatorname{Hom}(\underline\omega^{-1},\underline\omega)=\Gamma\underline\omega^2
\]
\uline{Questions}\qquad a) Quand $\alpha=0$ ?, i.e. $\bar\omega$ proportionnel : $\eta$ ??
b) Quand α est-il
b) Quand $\alpha$ est-il algébrique sur $\unsure{\mathbb Q|Q}$ ?
% archive sheet number: 66

